\documentclass[aspectratio=169, 9pt]{beamer}

% ── Thème et couleurs ────────────────────────────────────────
\usetheme{Madrid}
\usecolortheme{whale}
\usepackage[french]{babel}
\usepackage[utf8]{inputenc}
\usepackage[T1]{fontenc}
\usepackage{booktabs}
\usepackage{graphicx}
\usepackage{amsmath}
\usepackage{amssymb}
\usepackage{multirow}
\usepackage{xcolor}
\usepackage{tikz}
\usepackage{array}
\usepackage{colortbl}
\usepackage{adjustbox}
\usetikzlibrary{shapes, arrows, positioning,
                decorations.pathreplacing}

% ── Couleurs personnalisées ──────────────────────────────────
\definecolor{ugbblue}{RGB}{0, 51, 102}
\definecolor{ugbgold}{RGB}{204, 153, 0}
\definecolor{lightblue}{RGB}{173, 216, 230}
\definecolor{lightgold}{RGB}{255, 236, 179}
\definecolor{darkgreen}{RGB}{0, 100, 0}
\definecolor{lightgreen}{RGB}{144, 238, 144}

\setbeamercolor{structure}{fg=ugbblue}
\setbeamercolor{frametitle}{bg=ugbblue, fg=white}
\setbeamercolor{title}{fg=white}
\setbeamercolor{subtitle}{fg=ugbgold}
\setbeamercolor{author}{fg=white}
\setbeamercolor{institute}{fg=white}
\setbeamercolor{date}{fg=white}
\setbeamercolor{block title}{bg=ugbblue, fg=white}
\setbeamercolor{block body}{bg=lightblue!30}
\setbeamercolor{alerted text}{fg=ugbgold}

% ── Réduire les marges des blocks ────────────────────────────
\setlength{\leftmargini}{0.4em}
\setlength{\leftmarginii}{0.3em}

% ── Pied de page ─────────────────────────────────────────────
\setbeamertemplate{footline}{
  \leavevmode
  \hbox{
    \begin{beamercolorbox}[wd=.33\paperwidth,
      ht=2.25ex, dp=1ex, center]{author in head/foot}
      \usebeamerfont{author in head/foot}
      \insertshortauthor
    \end{beamercolorbox}
    \begin{beamercolorbox}[wd=.34\paperwidth,
      ht=2.25ex, dp=1ex, center]{title in head/foot}
      \usebeamerfont{title in head/foot}
      \insertshorttitle
    \end{beamercolorbox}
    \begin{beamercolorbox}[wd=.33\paperwidth,
      ht=2.25ex, dp=1ex, right]{date in head/foot}
      \usebeamerfont{date in head/foot}
      \insertframenumber{} / \inserttotalframenumber
      \hspace{2ex}
    \end{beamercolorbox}
  }
}
\setbeamertemplate{navigation symbols}{}

% ── Réduire l'espace vertical dans les listes ────────────────
\setbeamertemplate{itemize/enumerate body begin}
  {\small}
\setbeamertemplate{itemize/enumerate subbody begin}
  {\footnotesize}

% ── Informations du document ─────────────────────────────────
\title[Modélisation hybride --- Débits extrêmes]{
  Modélisation statistique des débits extrêmes\\
  dans le bassin du fleuve Sénégal\\
  basée sur une approche hybride de Machine Learning
}
\author[Modou \textsc{[Nom]}]{Modou \textsc{[Nom de famille]}}
\institute[UGB Saint-Louis]{
  Université Gaston Berger de Saint-Louis\\
  U.F.R. Sciences Appliquées et Technologie\\
  Filière : Mathématiques et Applications /
  Sciences des données\\
  \vspace{0.2cm}
  \textbf{Encadré par :} Professeur El Hadji DEME
}
\date{Année Universitaire 2024--2025}

% ============================================================
\begin{document}
% ============================================================

% ── SLIDE 1 — Page de titre ──────────────────────────────────
\begin{frame}
\titlepage
\end{frame}

% ── SLIDE 2 — Table des matières ─────────────────────────────
\begin{frame}{Table des Matières}
\tableofcontents
\end{frame}

% ============================================================
\section{Introduction Générale}
% ============================================================

% ── SLIDE 3 — Contexte et problématique ──────────────────────
\begin{frame}{Introduction Générale : Contexte et Problématique}
\begin{columns}[T]
\begin{column}{0.55\textwidth}
  \begin{block}{Le bassin du fleuve Sénégal}
  \small
    \begin{itemize}
      \item Système fluvial transfrontalier :
            Guinée, Mali, Mauritanie, Sénégal
      \item \textbf{13 millions} d'habitants concernés
      \item Régime contrôlé par la mousson
            ouest-africaine (juin--octobre)
      \item Barrages Manantali et Diama
    \end{itemize}
  \end{block}
  \vspace{0.1cm}
  \begin{alertblock}{Dynamiques critiques}
  \small
    \begin{itemize}
      \item Instabilité hydroclimatique croissante
      \item Intensification des crues :
            \textbf{+40\%} depuis 2000
      \item Inondations exceptionnelles 2024
            (Bakel -- Saint-Louis)
    \end{itemize}
  \end{alertblock}
\end{column}
\begin{column}{0.43\textwidth}
  \begin{figure}
    \includegraphics[width=\textwidth]{carte_bassin.png}
    \caption{\small Bassin du fleuve Sénégal (OMVS)}
  \end{figure}
\end{column}
\end{columns}
\end{frame}

% ── SLIDE 4 — Limites et objectifs ───────────────────────────
\begin{frame}{Introduction Générale : Limites et Objectifs}
\begin{columns}[T]
\begin{column}{0.48\textwidth}
  \begin{alertblock}{Limites des modèles classiques}
  \small
    \begin{itemize}
      \item GEV/GPD : hypothèse de
            \textbf{stationnarité} souvent violée
      \item ML seul : performances limitées
            sur séries courtes ($n=49$ ans)
      \item Dynamiques non linéaires
            mal capturées
    \end{itemize}
  \end{alertblock}
  \vspace{0.1cm}
  \begin{block}{Solution : approche hybride}
  \small
    \begin{itemize}
      \item EVT : cadre probabiliste rigoureux
      \item ML : apprentissage non linéaire
      \item \textbf{RF-GEV} : le meilleur
            des deux approches
    \end{itemize}
  \end{block}
\end{column}
\begin{column}{0.48\textwidth}
  \begin{block}{Objectifs de l'étude}
  \small
    \begin{enumerate}
      \item Analyser les structures temporelles
            des débits (1971--2019)
      \item Estimer les niveaux de retour
            par GEV et GPD
      \item Évaluer RF et LSTM
      \item \textbf{Développer le modèle
            hybride RF-GEV}
      \item Fournir des estimations pour
            la gestion des risques
    \end{enumerate}
  \end{block}
\end{column}
\end{columns}
\end{frame}

% ============================================================
\section{Cadre Hydrologique et Données}
% ============================================================

% ── SLIDE 5 — Zone d'étude et données ────────────────────────
\begin{frame}{Cadre Hydrologique : Zone d'Étude et Données}
\begin{columns}[T]
\begin{column}{0.48\textwidth}
  \begin{block}{Quatre stations hydrométriques}
  \small
    \begin{tabular}{ll}
      \toprule
      \textbf{Station} & \textbf{Position} \\
      \midrule
      Daka Saidou & Haut bassin \\
      Gourbassi   & Haut bassin \\
      Kidira      & Bassin moyen \\
      Oualia      & Haut bassin \\
      \bottomrule
    \end{tabular}
  \end{block}
  \vspace{0.1cm}
  \begin{block}{Données utilisées}
  \small
    \begin{itemize}
      \item Période : \textbf{1971--2019}
      \item \textbf{49} années hydrologiques
            (mai--avril)
      \item \textbf{17 897} obs. journalières
            par station
      \item Source : OMVS
      \item Imputation NA : moyenne
            climatologique mensuelle
    \end{itemize}
  \end{block}
\end{column}
\begin{column}{0.48\textwidth}
  \begin{figure}
    \includegraphics[width=\textwidth]{carte_bassin.png}
    \caption{\small Localisation des stations}
  \end{figure}
\end{column}
\end{columns}
\end{frame}

% ── SLIDE 6 — Statistiques descriptives ──────────────────────
\begin{frame}{Cadre Hydrologique : Exploration des Données}
\begin{columns}[T]
\begin{column}{0.5\textwidth}
  \begin{block}{Résumé statistique --- débits journaliers
               (m\textsuperscript{3}/s)}
    \resizebox{\textwidth}{!}{%
    \begin{tabular}{lcccc}
      \toprule
      \textbf{Station} & \textbf{Moy.} &
      \textbf{Skew.} & \textbf{Kurt.} &
      \textbf{CV (\%)} \\
      \midrule
      Daka Saidou & 202,7 & 1,616 & 4,822 & 139 \\
      Gourbassi   & 105,0 & 2,955 & 14,29 & 177 \\
      Kidira      & 122,6 & 3,532 & 19,39 & 197 \\
      Oualia      & 117,7 & 3,217 & 18,64 & 172 \\
      \bottomrule
    \end{tabular}}
  \end{block}
  \vspace{0.1cm}
  \begin{alertblock}{Caractéristiques clés}
  \small
    \begin{itemize}
      \item \textbf{Asymétrie positive} : 1,6 à 3,5
      \item \textbf{Leptokurticité} très prononcée
      \item Forte variabilité (CV > 100\%)
      \item Saisonnalité : juin--octobre
    \end{itemize}
  \end{alertblock}
\end{column}
\begin{column}{0.48\textwidth}
  \begin{figure}
    \includegraphics[width=\textwidth]{series_temporelles.png}
    \caption{\small Séries temporelles des débits}
  \end{figure}
\end{column}
\end{columns}
\end{frame}

% ── SLIDE 7 — Saisonnalité ───────────────────────────────────
\begin{frame}{Cadre Hydrologique : Saisonnalité et Prétraitement}
\begin{columns}[T]
\begin{column}{0.48\textwidth}
  \begin{figure}
    \includegraphics[width=\textwidth]{saisonnalite.png}
    \caption{\small Distribution mensuelle des débits}
  \end{figure}
\end{column}
\begin{column}{0.48\textwidth}
  \begin{block}{Valeurs manquantes traitées}
  \small
    \begin{tabular}{lcc}
      \toprule
      \textbf{Station} & \textbf{NA bruts} &
      \textbf{NA traités} \\
      \midrule
      Daka Saidou & 941   & 0 \\
      Gourbassi   & 945   & 0 \\
      Kidira      & 1 829 & 0 \\
      Oualia      & 1 621 & 0 \\
      \bottomrule
    \end{tabular}
  \end{block}
  \vspace{0.1cm}
  \begin{block}{Méthode}
  \small
    Imputation par \textbf{moyenne climatologique mensuelle}
    → 0 NA résiduel → 49 années complètes retenues
  \end{block}
  \vspace{0.1cm}
  \begin{alertblock}{Saisonnalité}
  \small
    Pics de crue entre \textbf{juin et octobre}
    sous l'effet de la mousson ouest-africaine
  \end{alertblock}
\end{column}
\end{columns}
\end{frame}

% ============================================================
\section{Théorie des Valeurs Extrêmes (TVE)}
% ============================================================

% ── SLIDE 8 — Concepts fondamentaux GEV ──────────────────────
\begin{frame}{TVE : Concepts Fondamentaux --- Loi GEV}
\begin{columns}[T]
\begin{column}{0.5\textwidth}
  \begin{block}{Modélisation des maxima : Loi GEV}
  \small
    Blocs annuels : $M_n = \max(X_1, \ldots, X_n)$
    \vspace{0.2cm}

    \textbf{Théorème de Gnedenko (1943)} :
    \begin{equation*}
      G_\xi(x) = \exp\!\left(-\!\left[1 +
      \xi\!\left(\frac{x-\mu}{\sigma}
      \right)\right]^{-1/\xi}\right)
    \end{equation*}
    \textbf{Trois paramètres} :
    $\mu$ (position), $\sigma$ (échelle),
    $\xi$ (forme)
  \end{block}
  \vspace{0.1cm}
  \begin{block}{Interprétation de $\xi$}
  \small
    \begin{tabular}{clc}
      \toprule
      \textbf{$\xi$} & \textbf{Loi} &
      \textbf{Queue} \\
      \midrule
      $> 0$ & Fréchet & \textcolor{red}{Lourde} \\
      $= 0$ & Gumbel  & \textcolor{orange}{Exponentielle} \\
      $< 0$ & Weibull & \textcolor{darkgreen}{Bornée} \\
      \bottomrule
    \end{tabular}
  \end{block}
\end{column}
\begin{column}{0.48\textwidth}
  \begin{block}{Méthode des L-moments}
  \small
    \textbf{Hosking (1990)} -- choix justifié :
    \begin{itemize}
      \item Robustes sur petits échantillons
            ($n=49 < 100$)
      \item Recommandés par la WMO
      \item Résistants aux valeurs aberrantes
      \item MLE divergeait ($\xi = +6$, absurde)
    \end{itemize}
  \end{block}
  \vspace{0.1cm}
  \begin{block}{Méthode POT --- Loi GPD}
  \small
    \begin{itemize}
      \item Seuil $u$ : 95\textsuperscript{e} percentile
      \item Déclustering : 7 jours
      \item $x_T = u + \frac{\sigma}{\xi}
            \left[(\lambda T)^\xi - 1\right]$
    \end{itemize}
  \end{block}
\end{column}
\end{columns}
\end{frame}

% ── SLIDE 9 — Structures temporelles ────────────────────────
\begin{frame}{TVE : Structures Temporelles et Stationnarité}
\begin{columns}[T]
\begin{column}{0.5\textwidth}
  \begin{block}{Test de Mann-Kendall}
  \small
    \begin{tabular}{lccc}
      \toprule
      \textbf{Station} & \textbf{$\tau$} &
      \textbf{p-val} & \textbf{Sen} \\
      \midrule
      Daka Saidou &
      \textcolor{red}{0,195} &
      \textcolor{red}{0,049} &
      5,80 \\
      Gourbassi & 0,188 & 0,058 & -- \\
      Kidira &
      \textcolor{red}{\textbf{0,304}} &
      \textcolor{red}{\textbf{0,002}} &
      \textbf{16,68} \\
      Oualia & 0,113 & 0,259 & -- \\
      \bottomrule
    \end{tabular}
  \end{block}
  \vspace{0.1cm}
  \begin{block}{Tests ADF et KPSS}
  \small
    \begin{itemize}
      \item \textbf{Kidira} :
            non-stationnarité \textbf{confirmée}\\
            ADF $p=0{,}207$ + KPSS $p=0{,}019$
      \item Autres stations : résultats ambigus
    \end{itemize}
  \end{block}
  \vspace{0.1cm}
  \begin{alertblock}{Implication}
  \small
    Kidira : \textbf{+15,27 m³/s/an}
    → GEV non stationnaire obligatoire
  \end{alertblock}
\end{column}
\begin{column}{0.48\textwidth}
  \begin{figure}
    \includegraphics[width=\textwidth]{tendance_sen.png}
    \caption{\small Tendance des maxima annuels ---
             pente de Sen par station}
  \end{figure}
\end{column}
\end{columns}
\end{frame}

% ── SLIDE 10 — GEV stationnaire ──────────────────────────────
\begin{frame}{TVE : Modélisation GEV Stationnaire}
\begin{columns}[T]
\begin{column}{0.52\textwidth}
  \begin{block}{Paramètres GEV (L-moments)}
  \small
    \begin{tabular}{lcccc}
      \toprule
      \textbf{Station} & \textbf{$\mu$} &
      \textbf{$\sigma$} & \textbf{$\xi$} &
      \textbf{AIC} \\
      \midrule
      Daka Saidou & 1\,014 & 204 & $-$0,001 & 699,3 \\
      Gourbassi   & 622    & 317 & $+$0,040  & 721,2 \\
      Oualia      & 538    & 314 & $-$0,207  & -- \\
      \bottomrule
    \end{tabular}
  \end{block}
  \vspace{0.1cm}
  \begin{block}{Niveaux de retour GEV stat. (m³/s)}
  \small
    \resizebox{\textwidth}{!}{%
    \begin{tabular}{lcccc}
      \toprule
      \textbf{Station} & \textbf{T=10} &
      \textbf{T=50} & \textbf{T=100} & \textbf{T=200} \\
      \midrule
      Daka Saidou & 1\,472 & 1\,808 & 1\,949 & 2\,091 \\
      Gourbassi   & 1\,369 & 1\,962 & 2\,225 & 2\,494 \\
      Oualia      & 1\,103 & 1\,379 & 1\,470 & 1\,549 \\
      \bottomrule
    \end{tabular}}
  \end{block}
\end{column}
\begin{column}{0.46\textwidth}
  \begin{figure}
    \includegraphics[width=\textwidth]{diagnostic_gev_daka_saidou.png}
    \caption{\small Diagnostics GEV --- Daka Saidou}
  \end{figure}
\end{column}
\end{columns}
\end{frame}

% ── SLIDE 11 — GEV non stationnaire ──────────────────────────
\begin{frame}{TVE : GEV Non Stationnaire --- Kidira}
\begin{columns}[T]
\begin{column}{0.48\textwidth}
  \begin{block}{Spécification NS}
  \small
    $\mu(t) = \mu_0 + \mu_1 \times t$
    \vspace{0.15cm}

    \textbf{Paramètres estimés} :
    \begin{itemize}
      \item $\mu_0 = 794{,}55$ m³/s
      \item $\mu_1 = \textbf{15{,}27}$ m³/s/an
      \item $\sigma = 370{,}90$ m³/s
      \item $\xi = 0{,}076$
    \end{itemize}
  \end{block}
  \vspace{0.1cm}
  \begin{block}{Test LRT}
  \small
    \begin{itemize}
      \item Stat. = 11,08 ; \textbf{p = 0,0009}
      \item $\Delta$AIC = $-$9,1 points
    \end{itemize}
  \end{block}
  \vspace{0.1cm}
  \begin{block}{Niveaux de retour NS (m³/s)}
  \small
    \begin{tabular}{lcc}
      \toprule
      \textbf{T} & \textbf{1971} & \textbf{2019} \\
      \midrule
      10 ans  & 1\,338 & 2\,071 \\
      100 ans & 2\,470 & \textbf{3\,202} \\
      200 ans & 2\,845 & 3\,578 \\
      \bottomrule
    \end{tabular}
  \end{block}
\end{column}
\begin{column}{0.48\textwidth}
  \begin{figure}
    \includegraphics[width=\textwidth]{diagnostic_gev_ns_kidira.png}
    \caption{\small Évolution de $\mu(t)$ --- Kidira}
  \end{figure}
\end{column}
\end{columns}
\end{frame}

% ── SLIDE 12 — GPD résultats ─────────────────────────────────
\begin{frame}{TVE : Résultats GPD --- Modélisation POT}
\begin{columns}[T]
\begin{column}{0.5\textwidth}
  \begin{block}{Paramètres GPD estimés (MLE)}
  \small
    \resizebox{\textwidth}{!}{%
    \begin{tabular}{lcccc}
      \toprule
      \textbf{Station} & \textbf{Seuil} &
      \textbf{$\lambda$} & \textbf{$\sigma$} &
      \textbf{$\xi$} \\
      \midrule
      Daka Saidou & 810,9 & 1,918 & 287,7 & $-$0,189 \\
      Gourbassi   & 483,2 & 1,388 & 320,3 & $-$0,054 \\
      Kidira      & 600,0 & 1,347 & 340,2 & $+$0,183 \\
      Oualia      & 527,1 & 0,959 & 320,6 & $+$0,196 \\
      \bottomrule
    \end{tabular}}
  \end{block}
  \vspace{0.1cm}
  \begin{block}{Niveaux de retour GPD (m³/s)}
  \small
    \resizebox{\textwidth}{!}{%
    \begin{tabular}{lcccc}
      \toprule
      \textbf{Station} & \textbf{T=10} &
      \textbf{T=50} & \textbf{T=100} & \textbf{T=200} \\
      \midrule
      Daka Saidou & 1\,462 & 1\,690 & 1\,769 & 1\,838 \\
      Gourbassi   & 1\,269 & 1\,697 & 1\,870 & 2\,037 \\
      Kidira      & 1\,733 & 2\,756 & 3\,298 & 3\,914 \\
      Oualia      & 1\,439 & 2\,384 & 2\,892 & 3\,474 \\
      \bottomrule
    \end{tabular}}
  \end{block}
\end{column}
\begin{column}{0.48\textwidth}
  \begin{figure}
    \includegraphics[width=\textwidth]{niveaux_retour_gpd.png}
    \caption{\small Courbes de niveaux de retour GPD}
  \end{figure}
  \begin{alertblock}{Constat}
  \small
    Kidira et Oualia ($\xi > 0$) :
    queues \textbf{lourdes} → risque élevé
    de crues exceptionnelles
  \end{alertblock}
\end{column}
\end{columns}
\end{frame}

% ── SLIDE 13 — Comparaison GEV vs GPD ────────────────────────
\begin{frame}{TVE : Comparaison GEV vs GPD}
\begin{block}{Niveaux de retour comparés (m\textsuperscript{3}/s)}
\centering
\resizebox{0.92\textwidth}{!}{%
\begin{tabular}{llcccc}
  \toprule
  \textbf{Station} & \textbf{Méthode} &
  \textbf{T=10} & \textbf{T=50} &
  \textbf{T=100} & \textbf{T=200} \\
  \midrule
  \multirow{2}{*}{Daka Saidou}
    & GEV stat. & 1\,472 & 1\,808 & 1\,949 & 2\,091 \\
    & GPD       & 1\,462 & 1\,690 & 1\,769 & 1\,838 \\
  \midrule
  \multirow{2}{*}{Gourbassi}
    & GEV stat. & 1\,369 & 1\,962 & 2\,225 & 2\,494 \\
    & GPD       & 1\,269 & 1\,697 & 1\,870 & 2\,037 \\
  \midrule
  \multirow{2}{*}{Kidira}
    & GEV NS (2019) & 2\,071 & 2\,845 &
      \textbf{3\,202} & 3\,578 \\
    & GPD           & 1\,733 & 2\,756 &
      3\,298 & 3\,914 \\
  \midrule
  \multirow{2}{*}{Oualia}
    & GEV stat. & 1\,103 & 1\,379 & 1\,470 & 1\,549 \\
    & GPD       & 1\,439 & 2\,384 & 2\,892 & 3\,474 \\
  \bottomrule
\end{tabular}}
\end{block}
\vspace{0.15cm}
\begin{columns}
\begin{column}{0.5\textwidth}
  \begin{alertblock}{Constats clés}
  \small
    \begin{itemize}
      \item GEV/GPD convergent pour T court
      \item Divergence forte à Oualia
            (GEV=1\,470 vs GPD=2\,892)
      \item Kidira : GEV NS indispensable
    \end{itemize}
  \end{alertblock}
\end{column}
\begin{column}{0.48\textwidth}
  \begin{block}{Conclusion EVT}
  \small
    Les deux approches sont
    \textbf{complémentaires} → justifie
    l'intégration dans le modèle hybride
  \end{block}
\end{column}
\end{columns}
\end{frame}

% ============================================================
\section{Machine Learning}
% ============================================================

% ── SLIDE 14 — Principes ML ──────────────────────────────────
\begin{frame}{Machine Learning : Principes et Méthodologie}
\begin{columns}[T]
\begin{column}{0.5\textwidth}
  \begin{block}{Algorithmes mobilisés}
  \small
    \begin{itemize}
      \item \textbf{Random Forest} :
            ensemble d'arbres réduisant la variance
      \item \textbf{LSTM} : réseau récurrent
            pour séries temporelles longues
      \item \textbf{Modèle naïf} :
            $\hat{Y}_t = Y_{t-1}$ (baseline)
    \end{itemize}
  \end{block}
  \vspace{0.1cm}
  \begin{block}{Optimisation}
  \small
    \textbf{Grid Search} + TimeSeriesSplit
    (5 folds) :
    \begin{itemize}
      \item Respecte l'ordre chronologique
      \item Évite toute fuite d'information
    \end{itemize}
  \end{block}
  \vspace{0.1cm}
  \begin{alertblock}{Split temporel}
  \small
    Train : 1971--2009 (39 ans)\\
    Test  : 2010--2019 (10 ans)
  \end{alertblock}
\end{column}
\begin{column}{0.48\textwidth}
  \begin{block}{Features construites (8 variables)}
  \small
    \begin{tabular}{ll}
      \toprule
      \textbf{Feature} & \textbf{Description} \\
      \midrule
      annee\_norm & Tendance temporelle \\
      annee²      & Tendance non linéaire \\
      lag1        & Débit $t-1$ \\
      lag2        & Débit $t-2$ \\
      lag3        & Débit $t-3$ \\
      moy\_3      & Moy. mobile 3 ans \\
      moy\_5      & Moy. mobile 5 ans \\
      std\_3      & Écart-type 3 ans \\
      \bottomrule
    \end{tabular}
  \end{block}
\end{column}
\end{columns}
\end{frame}

% ── SLIDE 15 — Random Forest résultats ───────────────────────
\begin{frame}{Machine Learning : Résultats Random Forest}
\begin{columns}[T]
\begin{column}{0.52\textwidth}
  \begin{block}{Hyperparamètres optimaux}
  \small
    \resizebox{\textwidth}{!}{%
    \begin{tabular}{lccc}
      \toprule
      \textbf{Station} & \textbf{N arbres} &
      \textbf{Prof. max} & \textbf{Min feuille} \\
      \midrule
      Daka Saidou & 300 & 3      & 2 \\
      Gourbassi   & 500 & 3      & 3 \\
      Kidira      & 100 & Aucune & 2 \\
      Oualia      & 100 & 5      & 2 \\
      \bottomrule
    \end{tabular}}
  \end{block}
  \vspace{0.1cm}
  \begin{block}{Performances Train / Test}
  \small
    \resizebox{\textwidth}{!}{%
    \begin{tabular}{lcccc}
      \toprule
      \textbf{Station} &
      \multicolumn{2}{c}{\textbf{Train}} &
      \multicolumn{2}{c}{\textbf{Test}} \\
      & RMSE & MAPE & RMSE & MAPE \\
      \midrule
      Daka Saidou & 135 & 10,9\% & 413 & 29,5\% \\
      Gourbassi   & 235 & 36,9\% & 550 & 46,3\% \\
      Kidira      & 254 & 30,4\% & 909 & 42,3\% \\
      Oualia      & 307 & 38,9\% & 603 & 105,4\% \\
      \bottomrule
    \end{tabular}}
  \end{block}
\end{column}
\begin{column}{0.46\textwidth}
  \begin{figure}
    \includegraphics[width=\textwidth]{rf_importance_variables.png}
    \caption{\small Importance des variables --- RF}
  \end{figure}
  \begin{alertblock}{Constat}
  \small
    \textbf{lag1} dominant → forte persistance
    temporelle des débits maximaux
  \end{alertblock}
\end{column}
\end{columns}
\end{frame}

% ── SLIDE 16 — LSTM architecture et résultats ────────────────
\begin{frame}{Machine Learning : LSTM --- Architecture et Résultats}
\begin{columns}[T]
\begin{column}{0.48\textwidth}
  \begin{block}{Architecture (32 033 paramètres)}
  \small
    \begin{tabular}{lc}
      \toprule
      \textbf{Couche} & \textbf{Paramètres} \\
      \midrule
      LSTM (64) + BatchNorm + Drop(0,3) & 18\,944 \\
      LSTM (32) + BatchNorm + Drop(0,2) & 12\,544 \\
      Dense (16, ReLU) + Drop(0,1)      & 528 \\
      Dense (1)                         & 17 \\
      \bottomrule
    \end{tabular}
  \end{block}
  \vspace{0.1cm}
  \begin{block}{Performances LSTM (Test)}
  \small
    \begin{tabular}{lcc}
      \toprule
      \textbf{Station} & \textbf{RMSE} & \textbf{MAPE} \\
      \midrule
      Daka Saidou & 627 & 46,6\% \\
      Gourbassi   & 589 & 47,8\% \\
      Kidira      & 890 & \textbf{36,4\%} \\
      Oualia      & 965 & 171,0\% \\
      \bottomrule
    \end{tabular}
  \end{block}
  \begin{alertblock}{Meilleur résultat}
  \small
    Kidira : MAPE \textbf{36,4\%} vs Naïf 61,3\%
    → \textbf{$-$42\%}
  \end{alertblock}
\end{column}
\begin{column}{0.48\textwidth}
  \begin{figure}
    \includegraphics[width=\textwidth]{lstm_predictions.png}
    \caption{\small LSTM --- observé vs prédit (2010--2019)}
  \end{figure}
\end{column}
\end{columns}
\end{frame}

% ── SLIDE 17 — Comparaison ML ─────────────────────────────────
\begin{frame}{Machine Learning : Comparaison des Performances}
\begin{block}{RMSE et MAPE --- Naïf vs RF vs LSTM (jeu de test 2010--2019)}
\centering
\small
\begin{tabular}{lcccccc}
  \toprule
  \textbf{Station} &
  \multicolumn{2}{c}{\textbf{Naïf}} &
  \multicolumn{2}{c}{\textbf{RF}} &
  \multicolumn{2}{c}{\textbf{LSTM}} \\
  & RMSE & MAPE & RMSE & MAPE & RMSE & MAPE \\
  \midrule
  Daka Saidou &
  \textbf{284} & \textbf{22\%} &
  413 & 29\% & 627 & 47\% \\
  Gourbassi &
  \textbf{427} & 60\% &
  550 & \textbf{46\%} & 589 & 48\% \\
  Kidira &
  \textbf{649} & 61\% &
  909 & 42\% & 890 & \textbf{36\%} \\
  Oualia &
  702 & 78\% &
  \textbf{603} & \textbf{105\%} &
  965 & 171\% \\
  \bottomrule
\end{tabular}
\end{block}
\vspace{0.15cm}
\begin{columns}
\begin{column}{0.5\textwidth}
  \begin{alertblock}{Constats clés}
  \small
    \begin{itemize}
      \item Naïf meilleur RMSE : \textbf{3/4 stations}
      \item LSTM meilleur MAPE : Kidira ($-$42\%)
      \item RF meilleur RMSE : Oualia
      \item Aucun ML ne domine systématiquement
    \end{itemize}
  \end{alertblock}
\end{column}
\begin{column}{0.48\textwidth}
  \begin{block}{Explication}
  \small
    \begin{itemize}
      \item $n = 49$ ans → trop court pour ML seul
      \item Pas de variables climatiques exogènes
      \item Surapprentissage partiel
      \item \textbf{→ Nécessité du modèle hybride}
    \end{itemize}
  \end{block}
\end{column}
\end{columns}
\end{frame}

% ============================================================
\section{Modèle Hybride RF-GEV}
% ============================================================

% ── SLIDE 18 — Architecture hybride ──────────────────────────
\begin{frame}{Modèle Hybride RF-GEV : Architecture}
\begin{columns}[T]
\begin{column}{0.52\textwidth}
  \begin{alertblock}{Motivation}
  \small
    \begin{itemize}
      \item ML seul : performances limitées
      \item GEV seul : stationnarité souvent violée
      \item \textbf{Hybride} : combine les deux
    \end{itemize}
  \end{alertblock}
  \vspace{0.1cm}
  \begin{block}{Spécification RF-GEV}
    \begin{equation*}
      \mu(t) = \beta_0 + \beta_1\hat{Q}_{RF}(t)
      + \beta_2\text{Année}(t)
    \end{equation*}
    \begin{equation*}
      \log\sigma(t) = \gamma_0 + \gamma_1\hat{Q}_{RF}(t)
    \end{equation*}
    \begin{equation*}
      \xi = \text{constant}
    \end{equation*}
  \end{block}
  \vspace{0.1cm}
  \begin{block}{Procédure hors échantillon}
  \small
    $\hat{Q}_{RF}(t)$ obtenu par
    \textbf{TimeSeriesSplit (5 folds)}\\
    → Garantit l'absence de fuite d'information
  \end{block}
\end{column}
\begin{column}{0.44\textwidth}
\begin{center}
\begin{tikzpicture}[
  node distance=0.35cm,
  box/.style={rectangle, rounded corners,
              draw=ugbblue, fill=lightblue!30,
              text width=2.8cm, align=center,
              font=\small, minimum height=0.65cm}
]
\node[box] (feat) {Features\\hydrologiques};
\node[box, below=of feat] (rf)
  {Random Forest\\(hors échantillon)};
\node[box, below=of rf,
      fill=lightgold, draw=ugbgold] (qrf)
  {$\hat{Q}_{RF}(t)$ covariable};
\node[box, below=of qrf] (gev)
  {GEV non stationnaire\\(MLE)};
\node[box, below=of gev,
      fill=lightgreen, draw=darkgreen]
  (nr) {Niveaux de retour\\hybrides};

\draw[->, thick, ugbblue] (feat) -- (rf);
\draw[->, thick, ugbgold] (rf)   -- (qrf);
\draw[->, thick, ugbblue] (qrf)  -- (gev);
\draw[->, thick, darkgreen] (gev) -- (nr);
\end{tikzpicture}
\end{center}
\end{column}
\end{columns}
\end{frame}

% ── SLIDE 19 — Paramètres estimés ────────────────────────────
\begin{frame}{Modèle Hybride : Paramètres Estimés}
\begin{columns}[T]
\begin{column}{0.55\textwidth}
  \begin{block}{Paramètres du modèle RF-GEV}
    \resizebox{\textwidth}{!}{%
    \begin{tabular}{lcccccc}
      \toprule
      \textbf{Station} &
      \textbf{$\beta_0$} &
      \textbf{$\beta_1$} &
      \textbf{$\beta_2$} &
      \textbf{$\gamma_0$} &
      \textbf{$\gamma_1$} &
      \textbf{$\xi$} \\
      \midrule
      Daka Saidou & 757,5 & $-$1,7 & $-$0,7 &
      6,21 & 0,19 & \textcolor{darkgreen}{$-$0,490} \\
      Gourbassi & 698,2 & $-$213,0 & 278,6 &
      5,59 & 0,36 & \textcolor{darkgreen}{$-$0,149} \\
      Kidira & 982,1 & $-$449,4 & 701,5 &
      5,81 & 0,20 & \textcolor{darkgreen}{$-$0,173} \\
      Oualia & 597,9 & 53,6 & $-$4,0 &
      5,76 & 0,08 & \textcolor{red}{$+$0,238} \\
      \bottomrule
    \end{tabular}}
  \end{block}
  \vspace{0.1cm}
  \begin{block}{Corrélation RF-OOS / Observé}
  \small
    \begin{tabular}{lc}
      \toprule
      \textbf{Station} & \textbf{Corrélation} \\
      \midrule
      Daka Saidou & 0,270 \\
      Gourbassi   & 0,253 \\
      Kidira      & \textbf{0,390} \\
      Oualia      & 0,126 \\
      \bottomrule
    \end{tabular}
  \end{block}
\end{column}
\begin{column}{0.43\textwidth}
  \begin{figure}
    \includegraphics[width=\textwidth]{hybride_parametres_temporels.png}
    \caption{\small Évolution de $\mu(t)$ et $\sigma(t)$}
  \end{figure}
  \begin{alertblock}{Interprétation $\xi$}
  \small
    $\xi < 0$ : Daka, Gourbassi, Kidira
    → queue \textbf{bornée}\\
    $\xi > 0$ : Oualia → queue \textbf{lourde}
  \end{alertblock}
\end{column}
\end{columns}
\end{frame}

% ── SLIDE 20 — Comparaison AIC ───────────────────────────────
\begin{frame}{Modèle Hybride : Comparaison AIC}
\begin{columns}[T]
\begin{column}{0.52\textwidth}
  \begin{block}{Gain AIC : GEV stat. vs Hybride RF-GEV}
  \small
    \begin{tabular}{lccc}
      \toprule
      \textbf{Station} &
      \textbf{GEV stat.} &
      \textbf{Hybride} &
      \textbf{$\Delta$AIC} \\
      \midrule
      Daka Saidou & 699,3 & 490,2 &
      \textcolor{darkgreen}{\textbf{$-$209,1}} \\
      Gourbassi & 721,2 & 479,8 &
      \textcolor{darkgreen}{\textbf{$-$241,4}} \\
      Kidira & -- & 493,6 & -- \\
      Oualia & -- & 505,0 & -- \\
      \bottomrule
    \end{tabular}
    \vspace{0.2cm}

    \footnotesize $|\Delta\text{AIC}| > 10$ =
    amélioration substantielle
  \end{block}
  \vspace{0.1cm}
  \begin{alertblock}{Conclusion}
  \small
    Gains de \textbf{209,1} et \textbf{241,4} points
    → Le modèle hybride améliore
    \textbf{significativement} l'ajustement GEV
    pour Daka Saidou et Gourbassi
  \end{alertblock}
\end{column}
\begin{column}{0.46\textwidth}
  \begin{figure}
    \includegraphics[width=\textwidth]{hybride_qqplot.png}
    \caption{\small QQ-plot résidus standardisés ---
             modèle hybride}
  \end{figure}
\end{column}
\end{columns}
\end{frame}

% ── SLIDE 21 — Niveaux de retour hybrides ────────────────────
\begin{frame}{Modèle Hybride : Niveaux de Retour (2019)}
\begin{block}{Niveaux de retour hybrides RF-GEV (m\textsuperscript{3}/s)}
\centering
\small
\begin{tabular}{lcccccc}
  \toprule
  \textbf{Station} &
  \textbf{$\mu$(2019)} &
  \textbf{$\sigma$(2019)} &
  \textbf{T=10} & \textbf{T=50} &
  \textbf{T=100} & \textbf{T=200} \\
  \midrule
  Daka Saidou & 756,5 & 481,4 &
  1\,413 & 1\,594 &
  \textcolor{ugbgold}{\textbf{1\,636}} & 1\,666 \\
  Gourbassi & 966,1 & 374,2 &
  1\,682 & 2\,074 &
  \textcolor{ugbgold}{\textbf{2\,212}} & 2\,337 \\
  Kidira & 1\,340,2 & 485,1 &
  2\,245 & 2\,717 &
  \textcolor{ugbgold}{\textbf{2\,879}} & 3\,023 \\
  Oualia & 594,2 & 317,8 &
  1\,541 & 2\,640 &
  \textcolor{ugbgold}{\textbf{3\,252}} & 3\,972 \\
  \bottomrule
\end{tabular}
\end{block}
\vspace{0.15cm}
\begin{columns}
\begin{column}{0.55\textwidth}
  \begin{figure}
    \includegraphics[width=\textwidth]{hybride_niveaux_retour.png}
    \caption{\small Courbes hybride vs GEV stationnaire}
  \end{figure}
\end{column}
\begin{column}{0.43\textwidth}
  \begin{alertblock}{Points clés}
  \small
    \begin{itemize}
      \item Estimations conditionnelles à \textbf{2019}
      \item Intègre l'état hydrologique récent
      \item Cohérent avec GEV NS (Kidira)
      \item Plus flexible que GEV stat.
    \end{itemize}
  \end{alertblock}
\end{column}
\end{columns}
\end{frame}

% ── SLIDE 22 — Comparaison complète tous modèles ─────────────
\begin{frame}{Modèle Hybride : Comparaison Complète}
\begin{block}{Niveaux de retour $T=100$ ans
             (m\textsuperscript{3}/s)}
\centering
\small
\begin{tabular}{lcccc}
  \toprule
  \textbf{Station} &
  \textbf{GEV stat.} &
  \textbf{GPD} &
  \textbf{GEV NS} &
  \textbf{Hybride RF-GEV} \\
  \midrule
  Daka Saidou & 1\,949 & 1\,769 & -- &
  \textcolor{ugbgold}{\textbf{1\,636}} \\
  Gourbassi & 2\,225 & 1\,870 & -- &
  \textcolor{ugbgold}{\textbf{2\,212}} \\
  Kidira & -- & 3\,298 & 3\,202 &
  \textcolor{ugbgold}{\textbf{2\,879}} \\
  Oualia & 1\,470 & 2\,892 & -- &
  \textcolor{ugbgold}{\textbf{3\,252}} \\
  \bottomrule
\end{tabular}
\end{block}
\vspace{0.15cm}
\begin{columns}
\begin{column}{0.55\textwidth}
  \begin{figure}
    \includegraphics[width=\textwidth]{comparaison_tous_modeles_nr.png}
    \caption{\small Comparaison des courbes de niveaux
             de retour par méthode}
  \end{figure}
\end{column}
\begin{column}{0.43\textwidth}
  \begin{alertblock}{Constats}
  \small
    \begin{itemize}
      \item GEV/GPD convergent pour T court
      \item Hybride Kidira :
            estimations \textbf{conservatrices}
      \item Oualia : hybride produit les
            estimations \textbf{les plus élevées}
      \item Hybride intègre la dynamique récente
    \end{itemize}
  \end{alertblock}
\end{column}
\end{columns}
\end{frame}

% ============================================================
\section{Résultats et Discussion}
% ============================================================

% ── SLIDE 23 — Synthèse résultats EVT ────────────────────────
\begin{frame}{Résultats et Discussion : Synthèse EVT}
\begin{columns}[T]
\begin{column}{0.5\textwidth}
  \begin{block}{Résultats GEV}
  \small
    \begin{itemize}
      \item Non-stationnarité confirmée à
            \textbf{Kidira} ($p=0{,}002$)
      \item GEV NS : $+15{,}27$ m³/s/an
      \item Niveaux centenaux :
            1\,949--3\,202 m³/s
      \item L-moments : estimateurs robustes
            sur $n=49$
    \end{itemize}
  \end{block}
  \vspace{0.1cm}
  \begin{block}{Résultats GPD}
  \small
    \begin{itemize}
      \item Niveaux centenaux :
            1\,769--3\,298 m³/s
      \item $\xi < 0$ : Daka, Gourbassi
            (queues bornées)
      \item $\xi > 0$ : Kidira, Oualia
            (queues lourdes)
      \item Forte divergence à Oualia
            vs GEV stat.
    \end{itemize}
  \end{block}
\end{column}
\begin{column}{0.48\textwidth}
  \begin{alertblock}{Résultats Hybride RF-GEV}
  \small
    \begin{itemize}
      \item \textbf{Gain AIC = $-$209,1}
            (Daka Saidou)
      \item \textbf{Gain AIC = $-$241,4}
            (Gourbassi)
      \item Niveaux de retour 2019 :
            \begin{itemize}\small
              \item Daka : 1\,636 m³/s
              \item Gourbassi : 2\,212
              \item Kidira : 2\,879
              \item Oualia : 3\,252
            \end{itemize}
      \item Cadre plus flexible et réaliste
    \end{itemize}
  \end{alertblock}
  \vspace{0.1cm}
  \begin{block}{Application}
  \small
    Base pour la gestion des risques
    d'inondation dans le bassin du
    fleuve Sénégal
  \end{block}
\end{column}
\end{columns}
\end{frame}

% ── SLIDE 24 — Synthèse résultats ML ─────────────────────────
\begin{frame}{Résultats et Discussion : Synthèse ML}
\begin{columns}[T]
\begin{column}{0.5\textwidth}
  \begin{figure}
    \includegraphics[width=\textwidth]{comparaison_metriques_ml.png}
    \caption{\small Comparaison RMSE et MAPE ---
             Naïf vs RF vs LSTM}
  \end{figure}
\end{column}
\begin{column}{0.48\textwidth}
  \begin{block}{Métriques EVT (AIC)}
  \small
    \begin{tabular}{lc}
      \toprule
      \textbf{Modèle} & \textbf{AIC} \\
      \midrule
      GEV stat. & 699,3 -- 756,2 \\
      GEV NS (Kidira) & 747,1 \\
      GPD & 658,8 -- 1\,220,8 \\
      \textbf{Hybride} & \textbf{479,8 -- 505,0} \\
      \bottomrule
    \end{tabular}
  \end{block}
  \vspace{0.1cm}
  \begin{alertblock}{Message clé}
  \small
    EVT et ML répondent à des
    \textbf{objectifs distincts et complémentaires} :
    \begin{itemize}
      \item EVT → niveaux de retour probabilistes
      \item ML → prédiction temporelle
      \item \textbf{Hybride → les deux}
    \end{itemize}
  \end{alertblock}
\end{column}
\end{columns}
\end{frame}

% ============================================================
\section{Conclusion \& Recommandations}
% ============================================================

% ── SLIDE 25 — Limites ───────────────────────────────────────
\begin{frame}{Limites de l'Étude}
\begin{columns}[T]
\begin{column}{0.5\textwidth}
  \begin{alertblock}{Limites méthodologiques}
    \begin{itemize}
      \item \textbf{Taille de l'échantillon}
            ($n = 49$ années) → contrainte
            pour le ML
      \item \textbf{Absence de variables
            climatiques} exogènes
            (précipitations, ENSO, NAO)
      \item \textbf{Surapprentissage ML}
            partiel observé
      \item Imputation NA par moyenne
            mensuelle → lissage léger
            de la variabilité
    \end{itemize}
  \end{alertblock}
\end{column}
\begin{column}{0.48\textwidth}
  \begin{alertblock}{Limites de comparaison}
    \begin{itemize}
      \item EVT et ML ont des objectifs
            \textbf{différents} :
            \begin{itemize}
              \item EVT → extrêmes rares
              \item ML → prédiction future
            \end{itemize}
      \item La comparaison directe
            n'est pas toujours pertinente
      \item AIC GEV non disponible pour
            Kidira NS et Oualia
            (modèles différents)
    \end{itemize}
  \end{alertblock}
  \vspace{0.1cm}
  \begin{block}{Perspective immédiate}
  \small
    Intégration de données pluviométriques
    pour enrichir les covariables du
    modèle hybride
  \end{block}
\end{column}
\end{columns}
\end{frame}

% ── SLIDE 26 — Perspectives ──────────────────────────────────
\begin{frame}{Perspectives de Recherche}
\begin{columns}[T]
\begin{column}{0.5\textwidth}
  \begin{block}{Court terme}
    \begin{itemize}
      \item Intégration de variables
            climatiques \textbf{ENSO/NAO}
            comme covariables
      \item Extension de la période
            d'étude \textbf{au-delà de 2020}
      \item Validation sur d'autres
            bassins sahéliens
    \end{itemize}
  \end{block}
  \vspace{0.2cm}
  \begin{block}{Moyen terme}
    \begin{itemize}
      \item Architecture hybride
            \textbf{LSTM-GEV}
      \item Modèle multi-sites
            (4 stations simultanément)
      \item Incertitude bootstrap
            sur les niveaux de retour
    \end{itemize}
  \end{block}
\end{column}
\begin{column}{0.48\textwidth}
  \begin{block}{Long terme -- Applications}
    \begin{itemize}
      \item \textbf{Dimensionnement} des
            ouvrages hydrauliques
      \item Système d'alerte précoce
            aux crues
      \item Planification de la gestion
            de l'eau par l'\textbf{OMVS}
      \item Modélisation des impacts
            du changement climatique
    \end{itemize}
  \end{block}
  \vspace{0.1cm}
  \begin{alertblock}{Impact attendu}
  \small
    Les niveaux de retour hybrides
    constituent une base précieuse
    pour la \textbf{réduction de la
    vulnérabilité} des populations
    riveraines du fleuve Sénégal
  \end{alertblock}
\end{column}
\end{columns}
\end{frame}

% ── SLIDE 27 — Conclusion ────────────────────────────────────
\begin{frame}{Conclusion Générale}
\begin{columns}[T]
\begin{column}{0.52\textwidth}
  \begin{block}{Bilan de l'étude}
    \begin{enumerate}
      \item \textbf{Non-stationnarité}
            confirmée à Kidira\\
            → GEV NS nécessaire et justifiée
      \item \textbf{GEV/GPD} : niveaux de
            retour robustes (1\,769--3\,298 m³/s)
      \item \textbf{ML} : performances
            limitées sur $n=49$ ans
      \item \textbf{Hybride RF-GEV} :
            gains AIC significatifs\\
            ($-$209,1 et $-$241,4 points)
      \item Cadre de modélisation
            \textbf{original et flexible}
    \end{enumerate}
  \end{block}
\end{column}
\begin{column}{0.46\textwidth}
  \begin{alertblock}{Contribution originale}
    Modèle hybride RF-GEV :
    \begin{itemize}
      \item Combine apprentissage non
            linéaire (RF) et cadre
            probabiliste des extrêmes (GEV)
      \item Validation temporelle hors
            échantillon rigoureuse
      \item Estimations conditionnelles
            à l'état hydrologique récent
    \end{itemize}
  \end{alertblock}
  \vspace{0.2cm}
  \begin{block}{Application directe}
  \small
    Outil opérationnel pour la gestion
    des risques d'inondation dans le
    bassin du fleuve Sénégal
  \end{block}
\end{column}
\end{columns}
\end{frame}

% ── SLIDE 28 — Recommandations ───────────────────────────────
\begin{frame}{Recommandations}
\begin{columns}[T]
\begin{column}{0.5\textwidth}
  \begin{block}{Pour la recherche}
    \begin{itemize}
      \item Enrichir les modèles avec des
            \textbf{données climatiques}
            (précipitations, ENSO, NAO)
      \item Explorer l'architecture
            \textbf{LSTM-GEV} hybride
      \item Étendre l'étude à d'autres
            bassins de la sous-région
      \item Quantifier l'incertitude
            par bootstrap paramétrique
    \end{itemize}
  \end{block}
\end{column}
\begin{column}{0.48\textwidth}
  \begin{block}{Pour la gestion des risques}
    \begin{itemize}
      \item Utiliser les niveaux de retour
            hybrides pour le
            \textbf{dimensionnement} des
            ouvrages hydrauliques
      \item Intégrer le modèle dans un
            \textbf{système d'alerte précoce}
            piloté par l'OMVS
      \item Mettre à jour les estimations
            avec les données post-2020
      \item Développer des cartes de
            \textbf{risques d'inondation}
            par période de retour
    \end{itemize}
  \end{block}
\end{column}
\end{columns}
\end{frame}

% ── SLIDE 29 — Merci ─────────────────────────────────────────
\begin{frame}{}
\begin{center}
  \vspace{0.5cm}
  {\Large\textcolor{ugbblue}{\textbf{Merci de votre attention !}}}

  \vspace{0.8cm}
  {\large\textcolor{ugbgold}{\textbf{Questions ?}}}

  \vspace{1cm}
  \begin{columns}
  \begin{column}{0.5\textwidth}
    \begin{block}{Travaux cités}
    \small
      \begin{itemize}
        \item Coles (2001), Hosking (1990)
        \item WMO (2009)
        \item Breiman (2001), Hochreiter (1997)
        \item Reichstein et al. (2019)
        \item Burnham \& Anderson (2002)
      \end{itemize}
    \end{block}
  \end{column}
  \begin{column}{0.45\textwidth}
    \begin{block}{Contact}
    \small
      \textbf{Modou [Nom de famille]}\\
      Université Gaston Berger\\
      Saint-Louis, Sénégal\\
      \vspace{0.2cm}
      Directeur : Pr. El Hadji DEME\\
      Co-enc. : Dr. Sadibou Aïdara
    \end{block}
  \end{column}
  \end{columns}
\end{center}
\end{frame}

% ============================================================
\end{document}
